% =====================================================================
%  Chapter 15 — Why More Data Can Make the Answer Worse
% =====================================================================

\chapter{Why More Data Can Make the Answer Worse}
\label{ch:more-data}

There is a feature of impact monitoring\idx{Monitoring} that unsettles people who meet it
for the first time, and it is worth naming plainly because it undermines the
simplest possible story about how the discipline works. A better orbit can
produce a \emph{higher} impact probability. A more accurate measurement can
make an object look more dangerous than it did before. And a careful
characterisation campaign can move an object \emph{out} of a keyhole and
leave the probability higher than it started.

None of this is a malfunction. All of it follows from the fact that the
impact probability is a conditional probability, and conditions renormalise.

\section{Where the probability actually comes from}

The impact probability is not a property of the object. It is a property of
the object \emph{given the observations}, and it is computed by propagating
the observation-conditioned distribution of the object's position to the
encounter epoch and asking what fraction of that distribution collides.

Two consequences follow immediately.

The first is that the probability is always conditional on something. It is
conditional on the orbital solution, on the arc, on the data reduction, on
the physical parameters. When any of those changes, the probability changes,
and it changes in whichever direction the new information points --- which
may be upward.

The second is that the probability is a density ratio, and a density ratio
moves when the denominator moves. If a refinement of the orbit removes
probability mass from the region that would have collided, the total
probability mass has to go somewhere, and the probability of collision rises.

\begin{seamnote}{The regression that alarms people}
The 2024~YR4 episode is a live instance of the general phenomenon, and the
NASA review of it describes the sequence of revisions to the object's
Earth-impact probability as the orbit was refined. The revisions did not
monotonically decrease the number. This is expected behaviour, and an
organisation that treats a rising number as a failure will make exactly the
wrong decisions.
\end{seamnote}

\section{Why more astrometry cannot fix a physical error}

There is a structural asymmetry here, and it is the reason the previous
chapter's arithmetic has the force it does.

The astrometric uncertainty enters the probability through \(\sigma_\zeta\),
the width of the along-track distribution. More astrometry --- more
observations, a longer arc, a better weighted solution --- reduces
\(\sigma_\zeta\), and the keyhole's probability falls as the distribution
narrows around a fixed keyhole position. The mechanism is well behaved and
monotone, and it is the reason the community's answer to a tail keyhole is
always ``observe more''.

The physical uncertainty enters the probability through the keyhole's
\emph{position}. A better measurement of the diameter or the density changes
where the object will be at the encounter epoch, and therefore changes
whether the keyhole lands on it. This mechanism is not monotone in either
direction. It can move the object into the keyhole or out of it.

The consequence for a programme is uncomfortable but plain. \textbf{There is
no quantity of astrometry that removes a physical error from an impact
probability.} The astrometric term can be driven towards zero; the physical
term can only be driven towards zero by measuring the physics.

\begin{table}[!htbp]
\centering
\small
\caption{The two terms, and what reduces each of them.}
\label{tab:terms}
\begin{tabularx}{\textwidth}{@{}P{1.28in}P{1.28in}Y@{}}
\toprule
\rowcolor{panel}
\textbf{Term} & \textbf{Reduced by} & \textbf{Behaviour} \\
\midrule
Astrometric & more observations, longer arc & monotone; the probability
falls \\
\rowcolor{panel}
Physical & measuring the size, the density, the spin state & not monotone; the
probability may rise or fall \\
\bottomrule
\end{tabularx}
\end{table}

\section{The three forms of the paradox}

The first is the refinement regression: a better orbit raises the
probability because the previous probability was an artefact of the
conditioning. The community's response is correct --- keep observing --- and
the number falls again once the solution is better still.

The second is the physical update: a measurement of the object's diameter or
density moves the keyhole centre and the probability moves with it, in
whichever direction. Here the community's response must be the opposite of
the first. Observing more will not help, because the error is not
astrometric. Measuring the physics is the only route, and it is the route
that is not funded.

The third is the keyhole exit. An object whose probability is high because it
sits in a keyhole has its probability \emph{fall} when the keyhole's centre
is refined away from it, and the fall can be dramatic, because the tail
magnifies the displacement in both directions. This is the good outcome, and
it can look like a mistake in the reporting.

\begin{table}[!htbp]
\centering
\small
\caption{The three behaviours, and what each one calls for.}
\label{tab:paradoxes}
\begin{tabularx}{\textwidth}{@{}P{1.32in}P{1.24in}Y@{}}
\toprule
\rowcolor{panel}
\textbf{Form} & \textbf{What it looks like} & \textbf{What is actually
required} \\
\midrule
Refinement regression\idx{Regression} & the probability rises as the orbit improves &
keep observing; this is the astrometric term \\
\rowcolor{panel}
Physical update & the probability moves when the size or density is
measured & measure more physics; observing will not help \\
Keyhole exit & the probability falls sharply as the centre is refined &
nothing further, but do not treat it as an error \\
\bottomrule
\end{tabularx}
\end{table}

\section{Why this is the argument for characterisation rather than for monitoring}

The monitoring\idx{Monitoring} machinery of the planetary-defence establishment is built on
the first and third of these behaviours and is good at them. It is not
built for the second, because the second requires an instrument, a target
list and an alert path, and none of the three exists in the shape the
situation needs. The documented position, from the 2025 NASA review of the
2024~YR4 episode, is that no automated system currently alerts physical
characterisation sites.

So the paradox is not only a scientific curiosity. It is the mechanism by
which the gap in Part I becomes dangerous rather than merely untidy. A
catalogue that contains 11,725 catalogued objects in the hazard class and
about 9,966 of them without any physical account contains, in statistical
form, 9,966 objects whose keyhole position is unknown. When one of those
objects turns out to be on a trajectory that passes near the Earth, the
monitoring\idx{Monitoring} machinery will do what it does: it will observe, refine, and
renormalise. And the probability it reports will depend on physical
assumptions that nobody has measured.

\begin{keytakeaways}
\item The impact probability is conditional, and a conditional probability
renormalises, so a better orbit can produce a higher probability.
\item The astrometric term is monotone and cheap to reduce; the physical
term is not monotone and no quantity of astrometry reduces it.
\item A physical update that raises the probability is not a malfunction,
and it is the case in which observing more is the wrong response.
\item Nine thousand nine hundred and sixty-six objects in the hazard class
carry an unknown keyhole position, and the machinery to measure it does not
exist.
\end{keytakeaways}
